\section{Implementation}
We implemented Heliod as a compile-time analysis inside of Abacus \cite{abacus}, a framework for typestate-based peripheral programming that uses Rust's procedural macro system. Heliod operates on a function's raw syntax tree, applying state inference and cost analysis rules to each visited node. Once completed, the code is passed through to the Rust type checker so that the code remains consisten with Rust's ownership and borrowing semantics. We primarily used Heliod on peripherals that are supported by Embassy, which is a Rust framework for developing embedded systems applications.

\subsection{Using Heliod}

To use Heliod, a programmer annotates a function, implementation block, or inline module with \texttt{\#[heliod(PeripheralCost)]}. Peripherals are declared using Abacus’s \texttt{process\_register\_block} machinery, which already provides state and substate declarations, stable and transient classifications, and register-operation permissions. Heliod reuses these as they are. The transition and downgrade lattice of Section 3 is shorthand for the same finite-state machine, so Heliod keeps no separate downgrade policy specification. Heliod adds a background cost to each state and a foreground cost to each operation. Costs may be a number of ticks or symbolic variables. Depending on how costs are provided, bounds are reported accordingly.  Furthermore, the programmer does not annotate peripheral states at use sites. A peripheral parameter without explicit annotations uses state inference to determine the set of states that the peripheral is in.

\subsection{Limitations \& Differences from the Formal Model}
The implementation and core calculus are closely related, however differ in a couple of important ways. First, due to how Embassy autogenerates register APIs using a PAC crate, Heliod is unable to be applied directly to the individual register accesses. The Abacus framework expects a fully annotated MMIO type to store all of the relevant information. While we could have reimplemented Embassy's autogenerated API to an annotated MMIO type, that would be difficult and potentially error-prone. Instead, Heliod is applied to internal helpers, which are small collections of register accesses and operations that compose together to form a well-defined API for the peripheral under analysis. These helpers declare their effect with \texttt{\#[heliod\_operation(peripheral, from, to, cost, background\_steps)]}, where \texttt{from} and \texttt{to} are states and \texttt{cost} can be either a symbolic cost variable or a number of ticks. the \texttt{background\_steps} is required as operations can last for more than one step, and specifying the number allows the analysis to propagate the background charge accurately. These are trusted contracts that Heliod uses to build up the cost analysis. 

For the actual analysis, the implementation is also able to reason about bounded loops by essentially performing loop unrolling. For unbounded loops, the implementation is able to reason about such loops where peripherals return to the same state they started in at the beginning of the loop, and there are no early returns. In such cases, Heliod will return a per-iteration bound to help aide developers. The implementation and model do not support more complex features of Rust, including arbitrary unsafe code, generics, closures, recursion and general trait implementations. However, Heliod allows for the use of the \texttt{heliod\_operation} macro as a contract for opaque functions which lie outside of Heliod's analysis, therefore allowing it to integrate with code it hasn't checked. 

\paragraph{Async Support}
Because Heliod needs to reason about hardware-software concurrency, it has a limited form of support for the Rust async/await programming model. Heliod is able to reason about asynchronous functions that are immediately awaited, along with Embassy's \texttt{OnDrop} and \texttt{with\_timeout()} fallback constructs for asynchronous programming. However anything beyond that is currently unsupported.